\documentclass[a4paper,12pt]{article}

\usepackage[T2A]{fontenc}
\usepackage[utf8]{inputenc}
\usepackage[russian]{babel}
\usepackage{amsmath,amssymb,mathtools}
\usepackage{geometry}
\usepackage{microtype}
\usepackage{icomma}
\usepackage{booktabs,tabularx,array}
\usepackage{enumitem}
\usepackage{fancyhdr}
\usepackage{tikz}
\usetikzlibrary{arrows.meta,positioning,shapes.geometric,calc}
\usepackage{hyperref}
\usepackage{xcolor}
\usepackage{setspace}
\usepackage{parskip}

\geometry{left=24mm,right=24mm,top=24mm,bottom=23mm,headheight=23pt}
\definecolor{accent}{HTML}{008080}
\definecolor{darkaccent}{HTML}{006B70}
\definecolor{ink}{HTML}{30343A}
\definecolor{quiet}{HTML}{686D72}
\hypersetup{hidelinks}
\setstretch{1.21}
\setlength{\parindent}{0pt}
\setlength{\parskip}{6pt}
\setlist[itemize]{leftmargin=6mm,itemsep=2pt,topsep=3pt}
\setlist[enumerate]{leftmargin=7mm,itemsep=4pt,topsep=3pt}
\renewcommand{\arraystretch}{1.17}

\pagestyle{fancy}
\fancyhf{}
\fancyhead[L]{\footnotesize\color{quiet}Дифференцированная схема кредитования}
\fancyhead[R]{\footnotesize\color{quiet}Анастасия Павлова}
\fancyfoot[C]{\footnotesize\thepage}
\renewcommand{\headrulewidth}{0.35pt}
\renewcommand{\headrule}{\hbox to\headwidth{\color{accent}\leaders\hrule height \headrulewidth\hfill}}

\newcommand{\mysection}[1]{%
  \vspace{8pt}\refstepcounter{section}%
  {\large\color{darkaccent}\thesection.\ #1}\par\vspace{3pt}%
  {\color{accent}\hrule height 0.45pt}\vspace{8pt}}
\newcommand{\term}[1]{\textcolor{darkaccent}{#1}}
\newcommand{\ex}[1]{\vspace{3pt}\textcolor{darkaccent}{\textbf{#1}}\par}

\begin{document}

\thispagestyle{empty}
\begin{center}
\vspace*{24mm}
{\large\color{darkaccent}МАТЕМАТИКА}

\vspace{25mm}
{\Huge Чек-лист}

\vspace{14mm}
{\LARGE\color{accent} Дифференцированная}\\[3pt]
{\LARGE\color{accent} схема кредитования}

\vspace{18mm}
{\large Подготовка к ЕГЭ по профильной математике}

\vfill
{\large Анастасия Павлова}

\vspace{11mm}
{\large 26 октября 2026 г.}

\vspace{21mm}
{\small Лёвин Артём Александрович}\\[2pt]
{\small эксклюзивно для Анастасии Павловой}

\vspace{16mm}
\begin{tikzpicture}
\draw[accent,line width=0.7pt] (0,0)
  .. controls (1.05,0.53) and (1.7,-0.32) .. (2.9,0)
  .. controls (4.05,0.31) and (4.6,-0.45) .. (5.8,0);
\fill[accent] (2.9,0) circle (1.5pt);
\end{tikzpicture}
\vspace{9mm}
\end{center}
\clearpage

\mysection{Как распознать схему}

\textbf{Дифференцированная схема кредитования:} основной долг
(\term{тело кредита}) уменьшается \textbf{на одну и ту же сумму}
в каждом платёжном периоде. Проценты начисляются на оставшуюся задолженность.

\textbf{Сигнал в задаче ЕГЭ:} «долг на определённую дату каждого года
(месяца) должен быть на одну и ту же сумму меньше долга на эту дату
предыдущего года».

\textbf{Сравнение:} в аннуитетной схеме общий платёж постоянен; в
дифференцированной схеме постоянна доля \emph{основного долга} в платеже.
При положительной ставке общий платёж постепенно уменьшается.

\mysection{Обозначения и пять колонок таблицы}

\begin{tabularx}{\linewidth}{@{}lX@{}}
\toprule
$S$ & первоначальная сумма кредита;\\
$n$ & число платёжных периодов (лет или месяцев);\\
$r$ & ставка \emph{за один период}, записанная десятичной дробью;\\
$d=S/n$ & постоянная часть платежа, погашающая тело кредита;\\
$A_k$ & полный платёж за $k$-й период.\\
\bottomrule
\end{tabularx}

Например, $10\%=0{,}10$. Если платят \emph{ежемесячно}, ставка $r$
должна относиться к месяцу; если \emph{ежегодно}~--- к году.

\ex{Пять колонок, общих для финансовых задач}
\begin{center}
\small
\begin{tabular}{@{}ccccc@{}}
\toprule
Начальный долг & Проценты & Общий долг & Платёж & Остаток\\
\midrule
$D_{k-1}$ & $rD_{k-1}$ & $D_{k-1}(1+r)$ &
$rD_{k-1}+d$ & $D_{k-1}-d$\\
\bottomrule
\end{tabular}
\end{center}

\mysection{Как устроен платёж}

Каждый платёж состоит ровно из двух частей:
\[
\boxed{\quad A_k=
  \underbrace{\frac{S}{n}}_{\text{погашение тела}}
 +\underbrace{r\left(S-\frac{(k-1)S}{n}\right)}_
  {\text{проценты на остаток}}\quad}
\]

После $k$ платежей долг равен
\[
D_k=S-\frac{kS}{n}.
\]
Поэтому $D_n=0$: за $n$ периодов тело кредита погашено полностью.

\ex{Пример 1. Таблица занятия: $S=120$, $r=10\%$, $n=3$}
Постоянная доля тела кредита: $d=120:3=40$.

\begin{center}
\begin{tabular}{@{}rrrrrr@{}}
\toprule
Год & Начало & Проценты & С процентами & Платёж & Остаток\\
\midrule
1 & 120 & 12 & 132 & $12+40=52$ & 80\\
2 &  80 &  8 &  88 & $8+40=48$  & 40\\
3 &  40 &  4 &  44 & $4+40=44$  &  0\\
\bottomrule
\end{tabular}
\end{center}

Общая сумма платежей $52+48+44=144$.
Переплата: $144-120=24=12+8+4$.

\mysection{Формулы, которые нужно знать}

В классической схеме с одной выплатой процентов и части тела за
каждый период, постоянной ставкой и равномерным погашением тела:

\[
\boxed{\text{ОСП}=S+\frac{Sr(n+1)}{2}}
\qquad
\boxed{\text{Переплата}=\frac{Sr(n+1)}{2}}
\]
\[
\boxed{\text{Процент переплаты}
   =\frac{r(n+1)}{2}\cdot100\%}
\]
\[
\boxed{A_1=Sr+\frac Sn}
\qquad
\boxed{A_n=\frac Sn\,r+\frac Sn}
\]

Здесь \term{ОСП}~--- \textbf{общая сумма платежей}, то есть все деньги,
отданные банку за весь срок. Первый платёж $A_1$~--- наибольший,
последний $A_n$~--- наименьший (при $r>0$ и $n>1$).

\textbf{Почему появляется $n+1$?} Проценты начисляются на
$S,\;S(n-1)/n,\;\ldots,\;S/n$. Это арифметическая прогрессия;
сумма её $n$ членов равна $S(n+1)/2$.

Соседние платежи различаются на постоянную сумму:
\[
A_k-A_{k+1}=\frac{Sr}{n}.
\]

\mysection{Три прикладных приёма из занятия}

\ex{Пример 2. Известен процент переплаты}
Кредит взят на $14$ месяцев, общая сумма выплат на $15\%$ больше
первоначальной суммы. Найдём \emph{месячную} ставку:
\[
\frac{r(14+1)}2\cdot100\%=15\%,
\qquad
r=\frac{2\cdot0{,}15}{15}=0{,}02.
\]
\term{Ответ:} $2\%$ в месяц.

\textbf{Важно:} $15\%$~--- переплата за весь срок, а $r$~--- ставка за период.

\ex{Пример 3. Неизвестен срок кредита}
Кредит $20$ млн рублей, ставка $15\%$ годовых, общая сумма выплат
$38$ млн рублей. Из формулы ОСП:
\[
38=20+\frac{20\cdot0{,}15(n+1)}2
=20+1{,}5(n+1).
\]
Отсюда $18=1{,}5(n+1)$, поэтому $n+1=12$.
\term{Ответ:} $11$ лет.

\ex{Пример 4. Неравенства для первого и последнего платежей}
Кредит $6$ млн рублей на $15$ лет. Наибольший платёж
\emph{не более} $1{,}9$ млн, наименьший~--- \emph{не менее} $0{,}5$ млн.
Постоянная доля тела равна $6:15=0{,}4$ млн.

\[
\begin{cases}
6r+0{,}4\le 1{,}9,\\
0{,}4r+0{,}4\ge 0{,}5
\end{cases}
\Longleftrightarrow
\begin{cases}
6r\le1{,}5,\\
0{,}4r\ge0{,}1
\end{cases}
\Longleftrightarrow
\begin{cases}
r\le0{,}25,\\
r\ge0{,}25.
\end{cases}
\]
Оба ограничения выполняются одновременно только при $r=0{,}25$.
\term{Ответ:} $25\%$ годовых.

\mysection{Четыре ошибки, которых нужно избежать}

\begin{enumerate}[label=\arabic*.]
\item \textbf{Ставка и переплата.} Процент переплаты за весь срок
и ставка за один период~--- разные величины. Отдельно подписывайте,
какая из них дана, а какая неизвестна.
\item \textbf{«Не более» и «не менее».} Записывайте неравенство
сначала словами: платёж \emph{не более} $b$ означает $A\le b$;
платёж \emph{не менее} $b$ означает $A\ge b$.
\item \textbf{Система неравенств.} Каждое неравенство решается
самостоятельно, а затем находится \emph{пересечение} множества решений.
Метод подстановки для систем уравнений здесь не применяется.
\item \textbf{Десятичные дроби.} Если ищете ставку $r$, сначала
найдите число $0{,}25$, затем переведите его в $25\%$.
Деление «процентного ответа» ещё на $100$ приводит к ошибке.
\end{enumerate}

\clearpage
\thispagestyle{fancy}
\begin{center}
{\Large\color{darkaccent}Алгоритм решения задачи о кредите}\\[3mm]
{\small Дифференцированная схема: равные доли основного долга}
\end{center}
\vspace{7mm}

\begin{center}
\begin{tikzpicture}[
  >=Stealth,
  node distance=10mm,
  action/.style={draw=accent,rounded corners=2pt,line width=.65pt,
    text width=100mm,minimum height=11mm,align=center,font=\normalsize},
  decision/.style={diamond,draw=accent,line width=.65pt,aspect=3.4,
    text width=50mm,inner sep=1.5mm,align=center,font=\small},
  terminal/.style={ellipse,draw=accent,line width=.65pt,
    text width=46mm,minimum height=10mm,align=center,font=\normalsize},
  arrow/.style={->,line width=.65pt,draw=darkaccent}
]
\node[terminal] (a) {Прочитайте условие};
\node[decision,below=of a] (b) {Долг уменьшается на одну и ту же сумму?};
\node[action,below=14mm of b] (c) {Запишите $S$, число периодов $n$, ставку $r$\\(если ставка неизвестна, оставьте букву $r$)};
\node[action,below=of c] (d) {Найдите постоянную долю тела: $d=S/n$};
\node[action,below=of d] (e) {Выберите нужную формулу: ОСП, переплата,\\первый или последний платёж};
\node[action,below=of e] (f) {Подставьте данные; при ограничениях запишите\\систему неравенств и найдите пересечение};
\node[action,below=of f] (g) {Проверьте ставку, единицы времени,\\направления знаков и итоговый ответ};
\node[terminal,below=of g] (h) {Ответ};
\node[action,right=12mm of b,text width=36mm,minimum height=13mm,font=\small] (other)
 {Уточните схему\\кредитования};
\draw[arrow] (a)--(b);
\draw[arrow] (b)--node[right,font=\scriptsize]{да}(c);
\draw[arrow] (b)--node[above,font=\scriptsize]{нет}(other);
\draw[arrow] (c)--(d);
\draw[arrow] (d)--(e);
\draw[arrow] (e)--(f);
\draw[arrow] (f)--(g);
\draw[arrow] (g)--(h);
\end{tikzpicture}
\end{center}

\clearpage
\begin{center}
{\LARGE\color{darkaccent}Тренировочный блок}\\[2mm]
{\small Путь А \,---\, 10 заданий от базовых к сложным}
\end{center}
\vspace{3mm}
Записывайте ставку как десятичную дробь при подстановке в формулу.
Если денежные суммы указаны в тысячах или миллионах рублей,
ответ сохраняйте в тех же единицах.

\textcolor{darkaccent}{\textbf{Базовый уровень}}

\begin{enumerate}[label=\textbf{\arabic*.},leftmargin=8mm,itemsep=7pt]
\item Кредит $90$ тыс.\ руб.\ выдан на $3$ года под $10\%$ годовых.
Найдите ежегодно погашаемую часть тела кредита и первый платёж.

\item Кредит $180$ тыс.\ руб.\ выдан на $3$ года под $10\%$
годовых. Составьте таблицу из пяти колонок, найдите ежегодные
платежи и общую сумму платежей.

\item Кредит $150$ тыс.\ руб.\ выдан на $5$ лет под $20\%$
годовых. Найдите платёж за третий год и долг после этого платежа.
\end{enumerate}

\textcolor{darkaccent}{\textbf{Применение формул}}

\begin{enumerate}[label=\textbf{\arabic*.},start=4,leftmargin=8mm,itemsep=7pt]
\item Кредит $480$ тыс.\ руб.\ выдан на $4$ года под $8\%$
годовых. Найдите переплату и общую сумму платежей.

\item Кредит $300$ тыс.\ руб.\ выдан на $5$ лет под $12\%$
годовых. Найдите наибольший и наименьший платежи и их разность.

\item Кредит $240$ тыс.\ руб.\ погашается равными долями тела за
$4$ года. Общая сумма платежей равна $300$ тыс.\ руб.
Найдите годовую процентную ставку.

\item Кредит взят на $8$ месяцев с равномерным уменьшением долга.
После его полного погашения сумма всех выплат оказалась на $18\%$
больше первоначальной суммы. Найдите месячную процентную ставку.
\end{enumerate}

\textcolor{darkaccent}{\textbf{Уровень выше среднего}}

\begin{enumerate}[label=\textbf{\arabic*.},start=8,leftmargin=8mm,itemsep=7pt]
\item В июле взят кредит $2$ млн руб.\ под $15\%$ годовых.
Ежегодно долг уменьшается на одну и ту же сумму.
После полного погашения общая сумма выплат равна $2{,}75$ млн руб.
На сколько лет был взят кредит?

\item Кредит $3$ млн руб.\ выдан на $6$ лет с равномерным
погашением тела. Наибольший платёж не превышает $0{,}85$ млн руб.,
наименьший составляет не менее $0{,}55$ млн руб.
Найдите все возможные годовые процентные ставки.

\item Кредит $5$ млн руб.\ выдан на $10$ лет с равномерным
погашением тела. Наибольший годовой платёж не более $1{,}25$ млн руб.,
наименьший~--- не менее $0{,}575$ млн руб.
Найдите годовую процентную ставку.
\end{enumerate}

\clearpage
\begin{center}
{\LARGE\color{darkaccent}Ответы}\\[3mm]
{\small Для самостоятельной проверки}
\end{center}

\begin{tabularx}{\linewidth}{@{}cX@{}}
\toprule
№ & Ответ\\
\midrule
1 & $30$ тыс.\ руб.; $39$ тыс.\ руб.\\
2 & Платежи $78$, $72$, $66$ тыс.\ руб.;
общая сумма $216$ тыс.\ руб.\\
3 & $48$ тыс.\ руб.; остаток $60$ тыс.\ руб.\\
4 & Переплата $96$ тыс.\ руб.; ОСП $576$ тыс.\ руб.\\
5 & $96$ тыс.\ руб.; $72$ тыс.\ руб.; разность $24$ тыс.\ руб.\\
6 & $10\%$ годовых.\\
7 & $4\%$ в месяц.\\
8 & $4$ года.\\
9 & От $10\%$ до $11\frac23\%$ годовых включительно.\\
10 & $15\%$ годовых.\\
\bottomrule
\end{tabularx}

\vspace{12mm}
{\large\color{darkaccent}Итоговая памятка}

\begin{tabularx}{\linewidth}{@{}p{49mm}X@{}}
\toprule
Сигнал в условии & Действие\\
\midrule
«На одну и ту же сумму меньше» &
Проверить дифференцированную схему; $d=S/n$.\\
Просят все выплаты & $\text{ОСП}=S+\frac{Sr(n+1)}2$.\\
Просят переплату & $\text{Переплата}=\frac{Sr(n+1)}2$.\\
Просят ставку при известной переплате &
Использовать $\frac{r(n+1)}2\cdot100\%$; отличать ставку от переплаты.\\
«Наибольший» / «наименьший» &
Применить формулу первого / последнего платежа.\\
«Не более» / «не менее» &
Записать $\le$ / $\ge$, решить систему неравенств.\\
\bottomrule
\end{tabularx}

\vfill
\begin{center}
\footnotesize\color{quiet}
Лёвин Артём Александрович\\
эксклюзивно для Анастасии Павловой
\end{center}

\end{document}
